\section{Implementation Details}

\subsection{Source Code and Directory Structure}

The system’s codebase is organized into a set of modular components, each encapsulating a distinct layer of functionality. 
This clear separation of concerns enhances maintainability, simplifies debugging, and reflects the architectural hierarchy of the system. 
Table~\ref{tab:source-code-structure} summarizes the purpose and responsibilities of the principal source files and directories.
% You might need to add \usepackage{longtable} to your preamble
\begin{longtable}{|p{4cm}|p{10cm}|}

\caption{Roles and responsibilities of major source code files.}
\label{tab:source-code-structure} \\

\hline
\textbf{File / Directory} & \textbf{Description} \\ 
\hline
\endfirsthead

\caption[]{-- continued from previous page} \\
\hline
\textbf{File / Directory} & \textbf{Description} \\ 
\hline
\endhead

\hline 
\multicolumn{2}{r}{{Continued on next page}} \\ 
\endfoot

\hline
\endlastfoot

% The table body starts here
\texttt{daemon/proxy.py} & Implements a multithreaded reverse proxy that serves as the single entry point for all client traffic. 
It listens on port~8080 and forwards incoming HTTP requests to the backend server according to the routing rules defined in \texttt{proxy.conf}. \\ 
\hline
\texttt{daemon/backend.py} & Provides the foundational low-level TCP server class responsible for accepting raw socket connections and delegating them to a higher-level request handler. 
This module abstracts generic networking behavior without embedding application-specific logic. \\ 
\hline
\texttt{daemon/weaprous.py} & A lightweight, custom-built web framework supplied as part of the assignment. 
It enables RESTful API development by allowing registration of URL routes and mapping them to corresponding handler functions (e.g., \texttt{@app.route}). \\ 
\hline
\texttt{daemon/p2p\_daemon.py} & The core module implementing peer-to-peer functionality. 
An instance of this class runs for each authenticated user, managing direct TCP connections with other peers. 
It supports the custom \texttt{CONNECT}/\texttt{ACCEPT} handshake and enforces the \texttt{PING}/\texttt{PONG} keepalive protocol. \\ 
\hline
\texttt{start\_sampleapp.py} & Serves as the main application server and the \textit{HTTP Bridge}. 
It defines all REST API endpoints for authentication, peer discovery, and P2P orchestration using the WeApRous framework. 
This module acts as the intermediary between the browser-based frontend and the background P2P daemons. \\ 
\hline
\texttt{static/js/chat.js} & The client-side JavaScript component responsible for rendering the user interface, handling interactions, and maintaining communication with the backend. 
It implements asynchronous long-polling mechanisms to receive real-time events and message updates from the server. \\ 
\hline
\texttt{daemon/tracker.py} & A central state management module that maintains a thread-safe, in-memory registry of all connected peers. 
It monitors peer status, coordinates connection updates, and enforces the 45-second heartbeat timeout for detecting inactive or disconnected clients. \\ 

\end{longtable}

\subsection{Core Module Implementations}

The system’s functionality is realized through three principal modules, each responsible for a distinct architectural layer: 
the low-level peer-to-peer (P2P) communication layer, the high-level HTTP application server, and the client-side web interface. 
This modular decomposition ensures a clear separation of responsibilities and facilitates scalability, testing, and maintenance.

\subsubsection{Backend -- P2P Daemon}

The decentralized communication layer is implemented in \texttt{p2p\_daemon.py}. 
An instance of the \texttt{P2PDaemon} class is spawned for each authenticated user, operating as an independent, multithreaded process that manages all direct TCP interactions between peers.

\paragraph{Connection Management:}
The daemon concurrently handles both inbound and outbound TCP connections. 
An internal \texttt{\_accept\_loop} thread continuously listens for incoming connection requests from other peers. 
When a request is detected, it is delegated to \texttt{\_handle\_incoming\_connection}, which performs the server-side handshake sequence (\texttt{CONNECT}~$\rightarrow$~\texttt{ACCEPT}/\texttt{REJECT}).  
For outbound connections, the \texttt{connect\_to\_peer()} method initiates the client-side handshake, establishing a verified communication channel with the target peer.

\paragraph{Message Loops and Timeout Detection:}
For each established P2P session, a dedicated \texttt{\_message\_loop} thread is created. 
This loop operates on a non-blocking socket configured with a 30-second idle timeout (\texttt{IDLE\_TIMEOUT}) to detect inactivity. 
Simultaneously, a \texttt{\_keepalive\_thread} periodically transmits a \texttt{PING} message every 10~seconds (\texttt{KEEPALIVE\_INTERVAL}). 
Together, these mechanisms enable robust detection of broken or unresponsive connections and ensure timely cleanup of inactive sessions.

\paragraph{Concurrency and Thread Safety:}
The module is heavily multithreaded to achieve full non-blocking operation. 
To prevent race conditions when managing the shared dictionary of active connections, all accesses are synchronized using \texttt{threading.RLock}. 
Incoming messages from peers are not processed inline but are instead forwarded to the main backend process through a thread-safe callback system, ensuring isolation between the networking and application layers.

\subsubsection{Backend -- HTTP Bridge}

The \texttt{start\_sampleapp.py} file serves as the main application server and acts as the critical "HTTP Bridge" that links the browser-based frontend with the P2P communication layer. 
It utilizes the \texttt{WeApRous} framework to define and manage all RESTful API endpoints. 
The module translates stateless HTTP requests into stateful operations performed by the corresponding \texttt{P2PDaemon} instance.

\paragraph{Key API Endpoints:}
\begin{itemize}
    \item \textbf{/api/p2p-connect:} Coordinates the P2P connection process. 
    It receives a request from the frontend, validates the target peer, and invokes the \texttt{connect\_to\_peer()} method on the appropriate daemon.
    \item \textbf{/api/p2p-send:} Relays chat messages. 
    It accepts message data from a frontend \texttt{POST} request and passes it to the \texttt{send\_message()} method of the sender’s daemon, which then transmits the message over the established TCP socket.
    \item \textbf{/api/p2p-receive:} Implements the server-side endpoint for message polling. 
    When a P2P daemon receives a \texttt{CHAT} message, it enqueues the message in a user-specific, thread-safe queue (\texttt{message\_queues}). 
    The frontend periodically polls this endpoint to retrieve queued messages and update the chat interface.
\end{itemize}

By abstracting the P2P subsystem behind a clean HTTP API, this module effectively overcomes browser-imposed restrictions on direct TCP socket access while preserving real-time responsiveness through asynchronous polling.

\subsubsection{Frontend -- Chat Web Application}

The client-side interface is implemented in \texttt{chat.js}, a lightweight JavaScript application responsible for rendering the user interface, processing user interactions, and maintaining communication with the backend.

\paragraph{Polling Mechanisms:}
The application simulates real-time communication through asynchronous polling routines executed at fixed intervals:
\begin{itemize}
    \item \textbf{pollPeerEvents():} Polls the \texttt{/api/broadcast-peer} endpoint every 3~seconds to receive updates on peers joining or leaving the network.
    \item \textbf{pollMessages():} Polls the \texttt{/api/p2p-receive} endpoint every 2~seconds to fetch newly received chat messages from the user’s server-side queue.
    \item \textbf{pollConnectionRequests()} and \textbf{pollConnectionResponses():} Manage the multi-step P2P connection handshake by polling connection request and response states.
\end{itemize}

\paragraph{UI and State Management:}
The script maintains client-side state through global variables such as \texttt{currentUser}, \texttt{activePeer}, and \texttt{connectedPeers}. 
It dynamically updates the Document Object Model (DOM) to reflect peer availability, connection status, and new message arrivals. 
Functions like \texttt{appendMessage()} and \texttt{updateConnectionStatus()} ensure that the interface remains synchronized with the backend and accurately reflects the state of the underlying network.

This frontend module, together with the backend bridge and P2P daemons, forms a cohesive end-to-end communication loop that supports real-time peer discovery, messaging, and fault-tolerant connectivity.


\subsection{Exception Handling and Debugging}

During development, several critical issues were identified and resolved to improve system stability and user experience. 
The debugging process provided significant insights into the challenges of asynchronous, multi-threaded, and multi-layered network applications. 
The three most notable bugs, as documented in the project’s \texttt{README.md}, are detailed below.

\subsubsection{Bug \#1: Indefinite Wait on Rejected Connection}

\paragraph{Problem Description:}
Initially, when User~A sent a connection request to User~B and User~B explicitly clicked ``Reject,’’ 
User~A’s interface would remain indefinitely stuck on a ``Waiting for response...’’ message. 
There was no mechanism to notify User~A that their request had been denied.

\paragraph{Root Cause Analysis:}
The system correctly implemented connection acceptance but lacked a complementary workflow for rejection handling. 
The backend did not record or forward rejection events, and the frontend had no endpoint to query for such status updates. 
As a result, the requesting client never received confirmation that the peer had refused the connection.

\paragraph{Solution:}
A complete rejection-handling workflow was introduced to address this issue:
\begin{itemize}
    \item \textbf{Backend:} A new list, \texttt{connection\_responses}, was added to \texttt{start\_sampleapp.py} to temporarily store both acceptance and rejection events. 
    A new API endpoint, \texttt{/api/p2p-get-responses}, was created to allow clients to retrieve these events.
    \item \textbf{Frontend:} A new polling function, \texttt{pollConnectionResponses()}, was implemented in \texttt{static/js/chat.js}, invoking the new endpoint every two seconds.
    \item \textbf{Logic:} When the frontend receives a rejection response, it immediately updates the user interface with a message (``Request rejected’’), closes the active chat window after a two-second delay, and redirects the user to the welcome screen.
\end{itemize}

\paragraph{Verification:}  
Testing was performed by having User~A initiate a connection while User~B clicked ``Reject.’’ 
User~A’s screen correctly displayed the rejection message, and the chat window closed as expected, confirming the fix’s effectiveness.

\subsubsection{Bug \#2: Chat Input Field Overlapped by Footer}

\paragraph{Problem Description:}
In the chat interface, when the number of messages exceeded the visible area, 
the message input field became obscured behind the application footer, 
preventing users from typing new messages once the chat history expanded.

\paragraph{Root Cause Analysis:}
This issue stemmed from incorrect CSS flexbox behavior. 
The parent container (\texttt{.chat-active}) did not properly constrain the height of its child elements, 
causing the message container to expand indefinitely and push the input field out of view.

\paragraph{Solution:}
The correction was implemented entirely in the CSS file \texttt{\seqsplit{static/css/chat.css}}:
\begin{itemize}
    \item The property \texttt{min-height: 0;} was applied to both \texttt{.chat-active} and \texttt{\seqsplit{.messages-container}} flex items to ensure they respected the parent container’s height constraints.
    \item The property \texttt{overflow: hidden;} was added to \texttt{.chat-active} to prevent unintended overflow beyond the viewport.
\end{itemize}

\noindent\textbf{Verification:}  
After applying the fix, the message container became independently scrollable, while the input field and ``Send’’ button remained fixed and visible at the bottom of the interface. 
Extensive testing confirmed that the layout behaved correctly regardless of message volume or window size.

\subsubsection{Bug \#3: Peer Not Notified When a Session Ends}

\paragraph{Problem Description:}
When User~A manually terminated a P2P session (e.g., by clicking an ``End Session’’ button), 
User~B received no notification. 
User~B’s chat window remained active, allowing further message input even though the connection had already been closed by the peer.

\paragraph{Root Cause Analysis:}
The initial implementation closed the TCP socket on User~A’s side without sending a termination notice to User~B. 
Consequently, the daemon on User~B’s end only detected the disconnection after the 30-second PING/PONG timeout, 
resulting in a delayed and confusing user experience.

\paragraph{Solution:}
A graceful disconnection workflow was added using the existing \texttt{CLOSE} message type in the P2P protocol:
\begin{itemize}
    \item \textbf{Protocol Enhancement:} The \texttt{CLOSE} message type was explicitly used to signal the end of a session.
    \item \textbf{Daemon Logic:} The P2P Daemon (\texttt{p2p\_daemon.py}) was modified to send a \texttt{CLOSE} message before closing its local socket.
    \item \textbf{Backend Bridge:} The backend server was updated to forward incoming \texttt{CLOSE} messages to the recipient’s message queue, just like standard \texttt{CHAT} messages.
    \item \textbf{Frontend Handling:} The \texttt{pollMessages()} function in \texttt{chat.js} now detects messages of type \texttt{CLOSE}. 
    Upon receiving one, it triggers \texttt{handleRemoteSessionEnd()}, which displays a notification ("Peer has ended the session"), disables the input field, and updates the interface to reflect the terminated connection.
\end{itemize}

\noindent\textbf{Verification:}  
After implementing this protocol enhancement, User~B’s interface immediately displayed a "Session ended by peer" message as soon as User~A terminated the connection. 
The input field was correctly disabled, and the UI state accurately reflected the end of the session, ensuring immediate and consistent feedback to both participants.

\subsection{Programming Standards and Compliance}

Adherence to professional programming standards and strict compliance with assignment constraints were maintained throughout the development of this project. 
The entire codebase was written to emphasize clarity, maintainability, and correctness, ensuring that the implementation remains consistent with both academic and industry best practices.

\subsubsection{Code Style and Documentation}
All Python source code strictly conforms to the \textbf{PEP~8} style guide, promoting uniform formatting, consistent naming conventions, and improved readability. 
Each major function, class, and module includes docstrings following the \textbf{PEP~257} convention, specifying their purpose, parameters, and return values. 
This disciplined approach makes the system largely self-documenting, facilitating ease of understanding and future extension.

\subsubsection{Standard Library Exclusivity}
In accordance with the assignment’s requirements, the entire project was implemented exclusively using the Python standard library. 
No third-party packages or external dependencies were installed or imported. 
The implementation primarily utilizes the following modules:
\begin{itemize}
    \item \texttt{socket} --- for low-level TCP/IP communication and connection management.
    \item \texttt{threading} --- for concurrent execution and thread synchronization.
    \item \texttt{argparse} --- for parsing command-line arguments and runtime configuration.
    \item \texttt{mimetypes} --- for dynamically identifying and serving static content with appropriate MIME types.
\end{itemize}
This design choice not only satisfies the assignment’s restrictions but also demonstrates mastery of Python’s built-in networking and concurrency capabilities.

\subsubsection{Framework-Free Implementation}
The system deliberately avoids high-level web frameworks such as Flask or Django. 
All HTTP server components, including request parsing, response generation, and routing logic, were developed manually using the provided \textit{WeApRous} micro-framework and core modules. 
This framework-free design ensures complete compliance with the assignment’s guidelines while highlighting the developers’ understanding of fundamental networking principles, protocol structures, and server architecture.

Overall, this strict adherence to programming and compliance standards underscores the project’s emphasis on foundational competence, reproducibility, and technical transparency.
